% AML (CSAI2017P) --- Theory Quiz 2 (T02) --- 8 sets, A--H, four pages each.
% Generated from the item store; do not edit this file directly.
% Released after the sitting. Worked solutions are issued separately.
\documentclass[10pt,a4paper]{article}
\usepackage[margin=16mm,top=12mm,bottom=14mm]{geometry}
\usepackage[T1]{fontenc}
\usepackage{lmodern}
\usepackage{enumitem}
\usepackage{fancyvrb}
\usepackage{amsmath}
\usepackage{amssymb}
\usepackage{array}
\usepackage{fancyhdr}
\usepackage{microtype}

\setlength{\parindent}{0pt}
\setlength{\emergencystretch}{3em}
\newcommand{\BL}{\rule[-1.5pt]{22mm}{0.4pt}}
\newcommand{\blk}[1]{\rule[-1.5pt]{#1}{0.4pt}}

\pagestyle{fancy}\fancyhf{}
\renewcommand{\headrulewidth}{0.4pt}
\renewcommand{\footrulewidth}{0.4pt}
\newcommand{\SetLetter}{}
\fancyhead[L]{\footnotesize CSAI2017P --- Applied Machine Learning (Theory) --- Theory Quiz 2}
\fancyhead[R]{\footnotesize\bfseries Set \SetLetter}
\fancyfoot[L]{\footnotesize Name:\,\blk{34mm} \quad SAP ID:\,\blk{28mm}}
\fancyfoot[R]{\footnotesize Set \SetLetter\ --- page \thepage\ of 4}

\DefineVerbatimEnvironment{code}{Verbatim}%
  {fontsize=\small,xleftmargin=5mm,baselinestretch=0.98}

\newlist{qA}{enumerate}{1}
\setlist[qA]{label=\textbf{A\arabic*.},leftmargin=8mm,topsep=4pt,itemsep=9pt,parsep=1pt}
\newlist{qB}{enumerate}{1}
\setlist[qB]{label=\textbf{B\arabic*.},leftmargin=8mm,topsep=5pt,itemsep=13pt,parsep=1pt}
\newlist{qC}{enumerate}{1}
\setlist[qC]{label=\textbf{C\arabic*.},leftmargin=8mm,topsep=5pt,itemsep=8pt,parsep=1pt}
\newlist{qD}{enumerate}{1}
\setlist[qD]{label=\textbf{D\arabic*.},leftmargin=8mm,topsep=3pt,itemsep=4pt,parsep=1pt}
\newlist{qE}{enumerate}{1}
\setlist[qE]{label=\textbf{E\arabic*.},leftmargin=8mm,topsep=3pt,itemsep=5pt,parsep=1pt}

\newcommand{\parthead}[2]{\vspace{2pt}\noindent\rule{\linewidth}{0.4pt}\par\vspace{2pt}%
  \noindent{\bfseries Part #1}\hfill{\small #2}\par\vspace{1pt}}

\makeatletter
\newcommand{\rules}[2]{\par\vspace{2mm}%
  \@tempcnta=0
  \loop\ifnum\@tempcnta<#1\relax
    \noindent\rule{\linewidth}{0.3pt}\par\vspace{#2}%
    \advance\@tempcnta by 1
  \repeat}
\makeatother

\begin{document}


\renewcommand{\SetLetter}{A}\setcounter{page}{1}


\begin{center}
{\large\bfseries CSAI2017P --- Applied Machine Learning (Theory)}\\[2pt]
{\large\bfseries Theory Quiz 2 (T02) \quad\textbullet\quad Set A}\\[3pt]
{\small Duration \textbf{40 minutes} \quad\textbullet\quad Maximum marks \textbf{20}
 \quad\textbullet\quad 15 questions on four pages}
\end{center}
\vspace{2pt}

% Column widths differ slightly from LQ01's: at LQ01's widths the row runs
% ~11.6pt past the margin and "Signature" does not fit its own cell.
\noindent\begin{tabular}{@{}p{1.4cm}p{5.0cm}p{1.9cm}p{3.5cm}p{1.3cm}p{2.4cm}@{}}
\textbf{Name} & \hrulefill & \textbf{SAP ID} & \hrulefill & \textbf{Batch} & \hrulefill\\[9pt]
\textbf{Date} & \hrulefill & \textbf{Signature} & \hrulefill & \textbf{Set} & \fbox{\rule{0pt}{11pt}\ A\ }\\
\end{tabular}

\vspace{4pt}
\noindent\fbox{\parbox{\dimexpr\linewidth-2\fboxsep-2\fboxrule}{\small
\textbf{Syllabus.} Prerequisite material --- probability, Python and elementary calculus --- together with
Unit I (Lecture 1) and Unit II (Lecture 2) as retained material: generalisation, cross-validation, and which
constant each loss fits; and, principally, \textbf{Unit III (Lecture 3)}: gradient descent and its update
rule, the learning rate and the stability threshold $\eta < 2/L$, batch, stochastic and mini-batch gradient
descent, and momentum.\\[3pt]
\textbf{Instructions.} Closed book. Write your answers on this paper, in the space provided. A
calculator is not needed --- the arithmetic is deliberately short. Answer every question --- there
is no negative marking. In Part C a reason is \textbf{optional}, but a correct reason can rescue a
wrong answer. In Parts D and E the marks are for the reasoning, not the vocabulary; an answer that
states what it cannot determine earns credit, one that asserts without a reason does not.\\[3pt]
\textbf{About the set letter.} Sets are handed out \textbf{at random}. Every set covers the same
topics at the same level, for the same marks, against the same marking scheme.}}
\vspace{3pt}


\parthead{A --- Multiple choice}{5 questions $\times$ 1 = 5 marks. Circle one letter.}

\begin{qA}
  \item A fair six-sided die is rolled once. $P(\text{the result is even})$ equals\\
    (a) $1/2$ \quad (b) $1/3$\\
    (c) $1/6$ \quad (d) $2/3$
  \item \texttt{len([w for w in range(10) if w \% 3 == 0])} evaluates to\\
    (a) 3 \quad (b) 4\\
    (c) 5 \quad (d) an error
  \item For $f(w) = (w-3)^2$, the derivative $f'(w)$ is\\
    (a) $(w-3)^2$ \quad (b) $2w-3$\\
    (c) $2(w-3)$ \quad (d) $w-3$
  \item Over constant predictions $c$, the quantity $\frac{1}{n}\sum_i (y_i - c)^2$ is minimised by the\\
    (a) median of the $y_i$ \quad (b) mode of the $y_i$\\
    (c) largest $y_i$ \quad (d) mean of the $y_i$
  \item For $f(w) = (w-3)^2$ the curvature is $f'' = 2$. Gradient descent on $f$ converges only while\\
    (a) $\eta < 1$ \quad (b) $\eta < 2$\\
    (c) $\eta < 0.5$ \quad (d) any $\eta > 0$
\end{qA}

\vspace{4pt}

\parthead{B --- Fill in the blanks}{4 $\times$ 1 = 4 marks. Write the missing word on the rule.}

\begin{qB}

  \item At a minimum of a smooth function the derivative equals \BL.

  \item Squared error is minimised by the mean; absolute error is minimised by the \BL.

  \item Writing $L$ for the largest curvature of $J$, gradient descent converges only while $\eta$ is smaller than \BL.

  \item With momentum $v \leftarrow \beta v + \nabla J$, the effective step size is $\eta$ divided by \BL.

\end{qB}

\vfill
\begin{center}\small
\begin{tabular}{|>{\centering\arraybackslash}m{1.9cm}|*{5}{>{\centering\arraybackslash}m{1.6cm}|}}
\hline
\rule{0pt}{12pt}\textbf{Question} & A1 & A2 & A3 & A4 & A5\\[2pt]\hline
\rule{0pt}{20pt}\textbf{Answer} & & & & & \\[7pt]\hline
\end{tabular}\\[2pt]
{\footnotesize Copy your circled Part A letters into this grid. The grid is what is marked.}
\end{center}

\newpage

\parthead{C --- True or false}{3 $\times$ 1 = 3 marks. Circle one. A reason is optional, but write one --- it can rescue a wrong answer.}

\begin{qC}

  \item A single train/test split on a small dataset gives a reliable estimate of how a model will perform on new data. \hfill \textbf{TRUE} \; / \; \textbf{FALSE}\\[1pt]
    {\small Reason:}\,\blk{140mm}\rules{1}{6.4mm}

  \item Doubling the learning rate always halves the number of iterations needed to converge. \hfill \textbf{TRUE} \; / \; \textbf{FALSE}\\[1pt]
    {\small Reason:}\,\blk{140mm}\rules{1}{6.4mm}

  \item On a well-conditioned problem, momentum can need \emph{more} iterations than plain gradient descent. \hfill \textbf{TRUE} \; / \; \textbf{FALSE}\\[1pt]
    {\small Reason:}\,\blk{140mm}\rules{1}{6.4mm}

\end{qC}

\vspace{4pt}

\parthead{D --- Short answer}{2 questions $\times$ 2 = 4 marks. Show how you decide, not just what you decide.}

\begin{qD}

  \item Minimise $f(w) = (w-4)^2$ by gradient descent from $w_0 = 0$ with $\eta = 0.25$, using $w \leftarrow w - \eta f'(w)$.
    \begin{enumerate}[label=(\alph*),leftmargin=7mm,itemsep=1pt,topsep=2pt]
      \item Compute $w_1$ and $w_2$, showing your working.
      \item State the error multiplier $\lvert 1 - \eta f''(w) \rvert$, and say whether this run converges monotonically, converges while oscillating, or diverges.
    \end{enumerate}\par\vspace{68mm}

\end{qD}

\newpage

\parthead{D --- Short answer, continued}{}

\begin{qD}[start=2]

  \item A model's training loss falls smoothly towards zero while its held-out error, measured at the same iterations, rises steadily. \textbf{(a)} Name what is happening. \textbf{(b)} State what you would change first, and why.\par\vspace{80mm}

\end{qD}

\vspace{4pt}

\parthead{E --- Reasoning}{1 question $\times$ 4 = 4 marks. You are marked on the reasoning you make visible.}

\begin{qE}

  \item A team fits a linear model by gradient descent on \textbf{standardised} features with $\eta = 0.05$;
it converges in a few hundred iterations. They then rebuild the same pipeline on the \textbf{raw} features
--- one of which is an amount in rupees running into the millions --- keep $\eta = 0.05$ unchanged, and the
loss becomes \texttt{nan} within three iterations.


    \vspace{2mm}
    Answer all three parts overleaf. Any diagnosis that fits the evidence and is argued for earns
    full credit --- naming the ``expected'' fault is not what is marked.
    \begin{enumerate}[label=(\roman*),leftmargin=7mm,itemsep=2pt,topsep=2pt]
      \item Name the most likely cause, and explain \emph{why} it makes the run diverge when the same
            $\eta$ was stable before. \hfill \textbf{[2]}
      \item State one concrete check that would confirm or rule out your diagnosis. \hfill \textbf{[1]}
      \item Say what you would expect to see if your diagnosis turned out to be wrong. \hfill \textbf{[1]}
    \end{enumerate}

\end{qE}

\newpage

\parthead{E --- Reasoning, answer space}{Label each part (i), (ii), (iii). A sketch, a table or a list is as acceptable as prose.}

\vfill

{\footnotesize\itshape Continue on a continuation sheet if you need more room; label anything you write there with its question number.}

\newpage

\renewcommand{\SetLetter}{B}\setcounter{page}{1}


\begin{center}
{\large\bfseries CSAI2017P --- Applied Machine Learning (Theory)}\\[2pt]
{\large\bfseries Theory Quiz 2 (T02) \quad\textbullet\quad Set B}\\[3pt]
{\small Duration \textbf{40 minutes} \quad\textbullet\quad Maximum marks \textbf{20}
 \quad\textbullet\quad 15 questions on four pages}
\end{center}
\vspace{2pt}

% Column widths differ slightly from LQ01's: at LQ01's widths the row runs
% ~11.6pt past the margin and "Signature" does not fit its own cell.
\noindent\begin{tabular}{@{}p{1.4cm}p{5.0cm}p{1.9cm}p{3.5cm}p{1.3cm}p{2.4cm}@{}}
\textbf{Name} & \hrulefill & \textbf{SAP ID} & \hrulefill & \textbf{Batch} & \hrulefill\\[9pt]
\textbf{Date} & \hrulefill & \textbf{Signature} & \hrulefill & \textbf{Set} & \fbox{\rule{0pt}{11pt}\ B\ }\\
\end{tabular}

\vspace{4pt}
\noindent\fbox{\parbox{\dimexpr\linewidth-2\fboxsep-2\fboxrule}{\small
\textbf{Syllabus.} Prerequisite material --- probability, Python and elementary calculus --- together with
Unit I (Lecture 1) and Unit II (Lecture 2) as retained material: generalisation, cross-validation, and which
constant each loss fits; and, principally, \textbf{Unit III (Lecture 3)}: gradient descent and its update
rule, the learning rate and the stability threshold $\eta < 2/L$, batch, stochastic and mini-batch gradient
descent, and momentum.\\[3pt]
\textbf{Instructions.} Closed book. Write your answers on this paper, in the space provided. A
calculator is not needed --- the arithmetic is deliberately short. Answer every question --- there
is no negative marking. In Part C a reason is \textbf{optional}, but a correct reason can rescue a
wrong answer. In Parts D and E the marks are for the reasoning, not the vocabulary; an answer that
states what it cannot determine earns credit, one that asserts without a reason does not.\\[3pt]
\textbf{About the set letter.} Sets are handed out \textbf{at random}. Every set covers the same
topics at the same level, for the same marks, against the same marking scheme.}}
\vspace{3pt}


\parthead{A --- Multiple choice}{5 questions $\times$ 1 = 5 marks. Circle one letter.}

\begin{qA}
  \item A bag holds $4$ red and $8$ blue balls. One is drawn at random. The probability that it is red is\\
    (a) $1/2$ \quad (b) $1/3$\\
    (c) $2/3$ \quad (d) $4$
  \item For \texttt{ws = [0, 3, 6, 9, 12]}, the expression \texttt{len(ws[1:4])} evaluates to\\
    (a) 2 \quad (b) 4\\
    (c) 3 \quad (d) 5
  \item The table gives $f(w) = (w-3)^2$ at $w = 2,\,3,\,4$ as $1,\,0,\,1$. The derivative $f'(w)$ is\\
    (a) $(w-3)^2$ \quad (b) $2w-3$\\
    (c) $w-3$ \quad (d) $2(w-3)$
  \item A column holds $2,\,3,\,4,\,5,\,96$. The single constant that minimises the squared error to these five values is\\
    (a) their mean, $22$ \quad (b) their median, $4$\\
    (c) their mode \quad (d) their largest value, $96$
  \item For $f(w) = (w-3)^2$ a table of $\lvert 1 - 2\eta \rvert$ reads $0.80,\ 0.00,\ 0.80,\ 1.10$ at $\eta = 0.10,\ 0.50,\ 0.90,\ 1.05$. The run that diverges is\\
    (a) $\eta = 0.90$ \quad (b) $\eta = 0.50$\\
    (c) $\eta = 1.05$ \quad (d) $\eta = 0.10$
\end{qA}

\vspace{4pt}

\parthead{B --- Fill in the blanks}{4 $\times$ 1 = 4 marks. Write the missing word on the rule.}

\begin{qB}

  \item A table of $f$ either side of its minimum is symmetric, and the derivative at the minimum itself equals \BL.

  \item For the values $2,\,3,\,4,\,5,\,96$ the mean is $22$ and the median is $4$. Absolute error is minimised by the \BL.

  \item With $L$ the largest curvature of $J$, gradient descent converges only while $\eta$ is smaller than \BL.

  \item A run with $\beta = 0.9$ takes steps about ten times the size of $\eta$, because the effective step is $\eta$ divided by \BL.

\end{qB}

\vfill
\begin{center}\small
\begin{tabular}{|>{\centering\arraybackslash}m{1.9cm}|*{5}{>{\centering\arraybackslash}m{1.6cm}|}}
\hline
\rule{0pt}{12pt}\textbf{Question} & A1 & A2 & A3 & A4 & A5\\[2pt]\hline
\rule{0pt}{20pt}\textbf{Answer} & & & & & \\[7pt]\hline
\end{tabular}\\[2pt]
{\footnotesize Copy your circled Part A letters into this grid. The grid is what is marked.}
\end{center}

\newpage

\parthead{C --- True or false}{3 $\times$ 1 = 3 marks. Circle one. A reason is optional, but write one --- it can rescue a wrong answer.}

\begin{qC}

  \item A table showing one train/test split's score on a small dataset is a reliable estimate of performance on new data. \hfill \textbf{TRUE} \; / \; \textbf{FALSE}\\[1pt]
    {\small Reason:}\,\blk{140mm}\rules{1}{6.4mm}

  \item A table of iteration counts must halve every time the learning rate doubles. \hfill \textbf{TRUE} \; / \; \textbf{FALSE}\\[1pt]
    {\small Reason:}\,\blk{140mm}\rules{1}{6.4mm}

  \item A table comparing $36$ iterations at $\beta = 0$ with $92$ at $\beta = 0.9$ on a well-conditioned problem is possible: momentum can need more iterations than plain gradient descent. \hfill \textbf{TRUE} \; / \; \textbf{FALSE}\\[1pt]
    {\small Reason:}\,\blk{140mm}\rules{1}{6.4mm}

\end{qC}

\vspace{4pt}

\parthead{D --- Short answer}{2 questions $\times$ 2 = 4 marks. Show how you decide, not just what you decide.}

\begin{qD}

  \item The table below is to be filled in for $f(w) = (w-5)^2$, from $w_0 = 0$, with $\eta = 0.6$ and $w \leftarrow w - \eta f'(w)$. \par\vspace{1mm}\hspace*{6mm}\begin{tabular}{@{}lll@{}}\hline\rule{0pt}{10pt}$k$ & $f'(w_{k-1})$ & $w_k$\\[1pt]\hline\rule{0pt}{10pt}$1$ & & \\ $2$ & & \\[1pt]\hline\end{tabular}
    \begin{enumerate}[label=(\alph*),leftmargin=7mm,itemsep=1pt,topsep=2pt]
      \item Compute $w_1$ and $w_2$, showing your working.
      \item State the error multiplier $\lvert 1 - \eta f''(w) \rvert$, and say whether this run converges monotonically, converges while oscillating, or diverges.
    \end{enumerate}\par\vspace{68mm}

\end{qD}

\newpage

\parthead{D --- Short answer, continued}{}

\begin{qD}[start=2]

  \item A results table records the training loss falling smoothly towards zero across the epochs while the held-out error, measured at the same epochs, rises steadily. \textbf{(a)} Name what is happening. \textbf{(b)} State what you would change first, and why.\par\vspace{80mm}

\end{qD}

\vspace{4pt}

\parthead{E --- Reasoning}{1 question $\times$ 4 = 4 marks. You are marked on the reasoning you make visible.}

\begin{qE}

  \item A team's results table records two runs of the same linear model, fitted by gradient descent with
$\eta = 0.05$ in both.

\vspace{2mm}
\begin{center}\small
\begin{tabular}{@{}lll@{}}
\hline
\rule{0pt}{11pt}\textbf{Run} & \textbf{Features} & \textbf{Outcome}\\[1pt]
\hline
\rule{0pt}{11pt}1 & standardised & converged, a few hundred iterations\\
2 & raw (one column is an amount in rupees, in the millions) & loss \texttt{nan} after 3 iterations\\[1pt]
\hline
\end{tabular}
\end{center}


    \vspace{2mm}
    Answer all three parts overleaf. Any diagnosis that fits the evidence and is argued for earns
    full credit --- naming the ``expected'' fault is not what is marked.
    \begin{enumerate}[label=(\roman*),leftmargin=7mm,itemsep=2pt,topsep=2pt]
      \item Name the most likely cause, and explain \emph{why} it makes the run diverge when the same
            $\eta$ was stable before. \hfill \textbf{[2]}
      \item State one concrete check that would confirm or rule out your diagnosis. \hfill \textbf{[1]}
      \item Say what you would expect to see if your diagnosis turned out to be wrong. \hfill \textbf{[1]}
    \end{enumerate}

\end{qE}

\newpage

\parthead{E --- Reasoning, answer space}{Label each part (i), (ii), (iii). A sketch, a table or a list is as acceptable as prose.}

\vfill

{\footnotesize\itshape Continue on a continuation sheet if you need more room; label anything you write there with its question number.}

\newpage

\renewcommand{\SetLetter}{C}\setcounter{page}{1}


\begin{center}
{\large\bfseries CSAI2017P --- Applied Machine Learning (Theory)}\\[2pt]
{\large\bfseries Theory Quiz 2 (T02) \quad\textbullet\quad Set C}\\[3pt]
{\small Duration \textbf{40 minutes} \quad\textbullet\quad Maximum marks \textbf{20}
 \quad\textbullet\quad 15 questions on four pages}
\end{center}
\vspace{2pt}

% Column widths differ slightly from LQ01's: at LQ01's widths the row runs
% ~11.6pt past the margin and "Signature" does not fit its own cell.
\noindent\begin{tabular}{@{}p{1.4cm}p{5.0cm}p{1.9cm}p{3.5cm}p{1.3cm}p{2.4cm}@{}}
\textbf{Name} & \hrulefill & \textbf{SAP ID} & \hrulefill & \textbf{Batch} & \hrulefill\\[9pt]
\textbf{Date} & \hrulefill & \textbf{Signature} & \hrulefill & \textbf{Set} & \fbox{\rule{0pt}{11pt}\ C\ }\\
\end{tabular}

\vspace{4pt}
\noindent\fbox{\parbox{\dimexpr\linewidth-2\fboxsep-2\fboxrule}{\small
\textbf{Syllabus.} Prerequisite material --- probability, Python and elementary calculus --- together with
Unit I (Lecture 1) and Unit II (Lecture 2) as retained material: generalisation, cross-validation, and which
constant each loss fits; and, principally, \textbf{Unit III (Lecture 3)}: gradient descent and its update
rule, the learning rate and the stability threshold $\eta < 2/L$, batch, stochastic and mini-batch gradient
descent, and momentum.\\[3pt]
\textbf{Instructions.} Closed book. Write your answers on this paper, in the space provided. A
calculator is not needed --- the arithmetic is deliberately short. Answer every question --- there
is no negative marking. In Part C a reason is \textbf{optional}, but a correct reason can rescue a
wrong answer. In Parts D and E the marks are for the reasoning, not the vocabulary; an answer that
states what it cannot determine earns credit, one that asserts without a reason does not.\\[3pt]
\textbf{About the set letter.} Sets are handed out \textbf{at random}. Every set covers the same
topics at the same level, for the same marks, against the same marking scheme.}}
\vspace{3pt}


\parthead{A --- Multiple choice}{5 questions $\times$ 1 = 5 marks. Circle one letter.}

\begin{qA}
  \item A spinner is divided into six equal sectors numbered $1$ to $6$. The probability that one spin lands on an even number is\\
    (a) $1/3$ \quad (b) $1/6$\\
    (c) $1/2$ \quad (d) $2/3$
  \item \texttt{len(set([0, 3, 3, 6, 9]))} evaluates to\\
    (a) 3 \quad (b) 5\\
    (c) an error \quad (d) 4
  \item A parabola $f(w) = (w-3)^2$ is drawn with its vertex at $w = 3$. The slope of the tangent at a point $w$ is\\
    (a) $2(w-3)$ \quad (b) $(w-3)^2$\\
    (c) $2w-3$ \quad (d) $w-3$
  \item A scatter of points is to be replaced by one horizontal line placed to make the total \emph{squared} vertical distance as small as possible. That line sits at the\\
    (a) median of the points \quad (b) mean of the points\\
    (c) mode of the points \quad (d) highest point
  \item On a plot of distance-to-minimum against iteration on a logarithmic scale, a straight line sloping \emph{upward} means\\
    (a) $\eta$ is well chosen \quad (b) $\eta$ is a little too small\\
    (c) the loss is not a quadratic \quad (d) $\eta$ is past the stability threshold
\end{qA}

\vspace{4pt}

\parthead{B --- Fill in the blanks}{4 $\times$ 1 = 4 marks. Write the missing word on the rule.}

\begin{qB}

  \item At the vertex of the parabola the tangent line is horizontal, so the derivative there equals \BL.

  \item One loss curve is a smooth parabola through the origin and is minimised by the mean; the other is a V with a corner at the origin and is minimised by the \BL.

  \item Writing $L$ for the largest curvature --- the steepest wall of the bowl --- gradient descent converges only while $\eta$ is smaller than \BL.

  \item On a plot of the weights $\beta^{\,j}$ against $j$, the areas sum to $1/(1-\beta)$; the effective step under momentum is therefore $\eta$ divided by \BL.

\end{qB}

\vfill
\begin{center}\small
\begin{tabular}{|>{\centering\arraybackslash}m{1.9cm}|*{5}{>{\centering\arraybackslash}m{1.6cm}|}}
\hline
\rule{0pt}{12pt}\textbf{Question} & A1 & A2 & A3 & A4 & A5\\[2pt]\hline
\rule{0pt}{20pt}\textbf{Answer} & & & & & \\[7pt]\hline
\end{tabular}\\[2pt]
{\footnotesize Copy your circled Part A letters into this grid. The grid is what is marked.}
\end{center}

\newpage

\parthead{C --- True or false}{3 $\times$ 1 = 3 marks. Circle one. A reason is optional, but write one --- it can rescue a wrong answer.}

\begin{qC}

  \item A single point on a performance chart, from one train/test split of a small dataset, is a reliable estimate of performance on new data. \hfill \textbf{TRUE} \; / \; \textbf{FALSE}\\[1pt]
    {\small Reason:}\,\blk{140mm}\rules{1}{6.4mm}

  \item If one loss curve is drawn for $\eta$ and another for $2\eta$, the second must always reach the floor in half as many iterations. \hfill \textbf{TRUE} \; / \; \textbf{FALSE}\\[1pt]
    {\small Reason:}\,\blk{140mm}\rules{1}{6.4mm}

  \item A plot showing momentum's trajectory taking a longer path than plain gradient descent on a circular bowl is possible: on a well-conditioned problem momentum can need more iterations. \hfill \textbf{TRUE} \; / \; \textbf{FALSE}\\[1pt]
    {\small Reason:}\,\blk{140mm}\rules{1}{6.4mm}

\end{qC}

\vspace{4pt}

\parthead{D --- Short answer}{2 questions $\times$ 2 = 4 marks. Show how you decide, not just what you decide.}

\begin{qD}

  \item A parabola $f(w) = (w-6)^2$ is sketched with its vertex at $w = 6$. A gradient-descent run starts at $w_0 = 0$ on the left-hand wall, with $\eta = 0.4$ and $w \leftarrow w - \eta f'(w)$.
    \begin{enumerate}[label=(\alph*),leftmargin=7mm,itemsep=1pt,topsep=2pt]
      \item Compute $w_1$ and $w_2$, showing your working.
      \item State the error multiplier $\lvert 1 - \eta f''(w) \rvert$, and say whether this run converges monotonically, converges while oscillating, or diverges.
    \end{enumerate}\par\vspace{68mm}

\end{qD}

\newpage

\parthead{D --- Short answer, continued}{}

\begin{qD}[start=2]

  \item Two curves are plotted against the epoch number: the training loss falls smoothly towards zero, while the held-out error falls at first and then turns upward and keeps rising. \textbf{(a)} Name what is happening from the turning point onward. \textbf{(b)} State what you would change first, and why.\par\vspace{80mm}

\end{qD}

\vspace{4pt}

\parthead{E --- Reasoning}{1 question $\times$ 4 = 4 marks. You are marked on the reasoning you make visible.}

\begin{qE}

  \item A team plots the loss against the iteration for a linear model fitted by gradient descent. On
\textbf{standardised} features with $\eta = 0.05$ the curve falls smoothly and flattens within a few hundred
iterations. They rebuild the same pipeline on the \textbf{raw} features --- one of which is an amount in
rupees running into the millions --- keep $\eta = 0.05$, and the new curve rises off the top of the axis and
becomes \texttt{nan} by the third iteration.


    \vspace{2mm}
    Answer all three parts overleaf. Any diagnosis that fits the evidence and is argued for earns
    full credit --- naming the ``expected'' fault is not what is marked.
    \begin{enumerate}[label=(\roman*),leftmargin=7mm,itemsep=2pt,topsep=2pt]
      \item Name the most likely cause, and explain \emph{why} it makes the run diverge when the same
            $\eta$ was stable before. \hfill \textbf{[2]}
      \item State one concrete check that would confirm or rule out your diagnosis. \hfill \textbf{[1]}
      \item Say what you would expect to see if your diagnosis turned out to be wrong. \hfill \textbf{[1]}
    \end{enumerate}

\end{qE}

\newpage

\parthead{E --- Reasoning, answer space}{Label each part (i), (ii), (iii). A sketch, a table or a list is as acceptable as prose.}

\vfill

{\footnotesize\itshape Continue on a continuation sheet if you need more room; label anything you write there with its question number.}

\newpage

\renewcommand{\SetLetter}{D}\setcounter{page}{1}


\begin{center}
{\large\bfseries CSAI2017P --- Applied Machine Learning (Theory)}\\[2pt]
{\large\bfseries Theory Quiz 2 (T02) \quad\textbullet\quad Set D}\\[3pt]
{\small Duration \textbf{40 minutes} \quad\textbullet\quad Maximum marks \textbf{20}
 \quad\textbullet\quad 15 questions on four pages}
\end{center}
\vspace{2pt}

% Column widths differ slightly from LQ01's: at LQ01's widths the row runs
% ~11.6pt past the margin and "Signature" does not fit its own cell.
\noindent\begin{tabular}{@{}p{1.4cm}p{5.0cm}p{1.9cm}p{3.5cm}p{1.3cm}p{2.4cm}@{}}
\textbf{Name} & \hrulefill & \textbf{SAP ID} & \hrulefill & \textbf{Batch} & \hrulefill\\[9pt]
\textbf{Date} & \hrulefill & \textbf{Signature} & \hrulefill & \textbf{Set} & \fbox{\rule{0pt}{11pt}\ D\ }\\
\end{tabular}

\vspace{4pt}
\noindent\fbox{\parbox{\dimexpr\linewidth-2\fboxsep-2\fboxrule}{\small
\textbf{Syllabus.} Prerequisite material --- probability, Python and elementary calculus --- together with
Unit I (Lecture 1) and Unit II (Lecture 2) as retained material: generalisation, cross-validation, and which
constant each loss fits; and, principally, \textbf{Unit III (Lecture 3)}: gradient descent and its update
rule, the learning rate and the stability threshold $\eta < 2/L$, batch, stochastic and mini-batch gradient
descent, and momentum.\\[3pt]
\textbf{Instructions.} Closed book. Write your answers on this paper, in the space provided. A
calculator is not needed --- the arithmetic is deliberately short. Answer every question --- there
is no negative marking. In Part C a reason is \textbf{optional}, but a correct reason can rescue a
wrong answer. In Parts D and E the marks are for the reasoning, not the vocabulary; an answer that
states what it cannot determine earns credit, one that asserts without a reason does not.\\[3pt]
\textbf{About the set letter.} Sets are handed out \textbf{at random}. Every set covers the same
topics at the same level, for the same marks, against the same marking scheme.}}
\vspace{3pt}


\parthead{A --- Multiple choice}{5 questions $\times$ 1 = 5 marks. Circle one letter.}

\begin{qA}
  \item A quality inspector picks one component at random from a tray of $12$, of which $4$ are marked for rework. The probability of picking a marked one is\\
    (a) $1/2$ \quad (b) $2/3$\\
    (c) $4$ \quad (d) $1/3$
  \item A shop's opening hours are stored as \texttt{\{'mon': 9, 'tue': 9, 'mon': 10\}}. Its \texttt{len} is\\
    (a) 2 \quad (b) 1\\
    (c) 3 \quad (d) an error
  \item A cost rises as the square of how far a setting sits from its ideal value $3$, so $f(w) = (w-3)^2$. The rate at which the cost changes with $w$ is\\
    (a) $(w-3)^2$ \quad (b) $2(w-3)$\\
    (c) $2w-3$ \quad (d) $w-3$
  \item A ward records stays of $2,\,3,\,4,\,5$ and $96$ days. A single number is wanted that minimises the total \emph{squared} error to all five. It is\\
    (a) the median, $4$ days \quad (b) the mode\\
    (c) the mean, $22$ days \quad (d) the longest stay, $96$ days
  \item An engineer's training run explodes to \texttt{nan} within a few iterations. Before anything else, the first suspect is that\\
    (a) the dataset is too small \quad (b) $\eta$ is past the stability threshold\\
    (c) the model has too few parameters \quad (d) the loss function is the wrong one
\end{qA}

\vspace{4pt}

\parthead{B --- Fill in the blanks}{4 $\times$ 1 = 4 marks. Write the missing word on the rule.}

\begin{qB}

  \item A machine is at its cheapest setting when the cost stops changing with the setting --- that is, when the derivative equals \BL.

  \item A delivery firm that wants the typical error small rather than the squared error small should fit the \BL\ of the times rather than their mean.

  \item An engineer who knows the largest curvature $L$ of the cost surface knows that the run converges only while $\eta$ is smaller than \BL.

  \item A colleague who sets $\beta = 0.9$ finds the model taking steps about ten times as large, because the effective step is $\eta$ divided by \BL.

\end{qB}

\vfill
\begin{center}\small
\begin{tabular}{|>{\centering\arraybackslash}m{1.9cm}|*{5}{>{\centering\arraybackslash}m{1.6cm}|}}
\hline
\rule{0pt}{12pt}\textbf{Question} & A1 & A2 & A3 & A4 & A5\\[2pt]\hline
\rule{0pt}{20pt}\textbf{Answer} & & & & & \\[7pt]\hline
\end{tabular}\\[2pt]
{\footnotesize Copy your circled Part A letters into this grid. The grid is what is marked.}
\end{center}

\newpage

\parthead{C --- True or false}{3 $\times$ 1 = 3 marks. Circle one. A reason is optional, but write one --- it can rescue a wrong answer.}

\begin{qC}

  \item A consultant who reports one train/test split's score on a small dataset has given a reliable estimate of how the model will do in production. \hfill \textbf{TRUE} \; / \; \textbf{FALSE}\\[1pt]
    {\small Reason:}\,\blk{140mm}\rules{1}{6.4mm}

  \item An engineer whose run takes $200$ iterations can be sure that doubling the learning rate will bring it down to $100$. \hfill \textbf{TRUE} \; / \; \textbf{FALSE}\\[1pt]
    {\small Reason:}\,\blk{140mm}\rules{1}{6.4mm}

  \item An engineer who adds momentum to a well-conditioned problem may find the run needs \emph{more} iterations than before. \hfill \textbf{TRUE} \; / \; \textbf{FALSE}\\[1pt]
    {\small Reason:}\,\blk{140mm}\rules{1}{6.4mm}

\end{qC}

\vspace{4pt}

\parthead{D --- Short answer}{2 questions $\times$ 2 = 4 marks. Show how you decide, not just what you decide.}

\begin{qD}

  \item A machine's cost is $f(w) = (w-4)^2$, where $w$ is a dial setting and the ideal setting is $4$. An engineer tunes the dial by gradient descent from $w_0 = 0$, using $w \leftarrow w - \eta f'(w)$ with $\eta = 0.75$.
    \begin{enumerate}[label=(\alph*),leftmargin=7mm,itemsep=1pt,topsep=2pt]
      \item Compute $w_1$ and $w_2$, showing your working.
      \item State the error multiplier $\lvert 1 - \eta f''(w) \rvert$, and say whether this run converges monotonically, converges while oscillating, or diverges.
    \end{enumerate}\par\vspace{68mm}

\end{qD}

\newpage

\parthead{D --- Short answer, continued}{}

\begin{qD}[start=2]

  \item An analyst reports that her churn model's training loss has fallen smoothly towards zero over the epochs, while the error on the customers held back for testing has risen steadily over the same epochs. \textbf{(a)} Name what is happening. \textbf{(b)} State what you would change first, and why.\par\vspace{80mm}

\end{qD}

\vspace{4pt}

\parthead{E --- Reasoning}{1 question $\times$ 4 = 4 marks. You are marked on the reasoning you make visible.}

\begin{qE}

  \item An engineer fits a linear model by gradient descent on \textbf{standardised} features with
$\eta = 0.05$, and it converges in a few hundred iterations. A colleague then rebuilds the pipeline from the
\textbf{raw} customer table --- where one column is a lifetime spend in rupees, running into the millions ---
keeps $\eta = 0.05$ because ``it worked last time'', and the loss becomes \texttt{nan} within three
iterations.


    \vspace{2mm}
    Answer all three parts overleaf. Any diagnosis that fits the evidence and is argued for earns
    full credit --- naming the ``expected'' fault is not what is marked.
    \begin{enumerate}[label=(\roman*),leftmargin=7mm,itemsep=2pt,topsep=2pt]
      \item Name the most likely cause, and explain \emph{why} it makes the run diverge when the same
            $\eta$ was stable before. \hfill \textbf{[2]}
      \item State one concrete check that would confirm or rule out your diagnosis. \hfill \textbf{[1]}
      \item Say what you would expect to see if your diagnosis turned out to be wrong. \hfill \textbf{[1]}
    \end{enumerate}

\end{qE}

\newpage

\parthead{E --- Reasoning, answer space}{Label each part (i), (ii), (iii). A sketch, a table or a list is as acceptable as prose.}

\vfill

{\footnotesize\itshape Continue on a continuation sheet if you need more room; label anything you write there with its question number.}

\newpage

\renewcommand{\SetLetter}{E}\setcounter{page}{1}


\begin{center}
{\large\bfseries CSAI2017P --- Applied Machine Learning (Theory)}\\[2pt]
{\large\bfseries Theory Quiz 2 (T02) \quad\textbullet\quad Set E}\\[3pt]
{\small Duration \textbf{40 minutes} \quad\textbullet\quad Maximum marks \textbf{20}
 \quad\textbullet\quad 15 questions on four pages}
\end{center}
\vspace{2pt}

% Column widths differ slightly from LQ01's: at LQ01's widths the row runs
% ~11.6pt past the margin and "Signature" does not fit its own cell.
\noindent\begin{tabular}{@{}p{1.4cm}p{5.0cm}p{1.9cm}p{3.5cm}p{1.3cm}p{2.4cm}@{}}
\textbf{Name} & \hrulefill & \textbf{SAP ID} & \hrulefill & \textbf{Batch} & \hrulefill\\[9pt]
\textbf{Date} & \hrulefill & \textbf{Signature} & \hrulefill & \textbf{Set} & \fbox{\rule{0pt}{11pt}\ E\ }\\
\end{tabular}

\vspace{4pt}
\noindent\fbox{\parbox{\dimexpr\linewidth-2\fboxsep-2\fboxrule}{\small
\textbf{Syllabus.} Prerequisite material --- probability, Python and elementary calculus --- together with
Unit I (Lecture 1) and Unit II (Lecture 2) as retained material: generalisation, cross-validation, and which
constant each loss fits; and, principally, \textbf{Unit III (Lecture 3)}: gradient descent and its update
rule, the learning rate and the stability threshold $\eta < 2/L$, batch, stochastic and mini-batch gradient
descent, and momentum.\\[3pt]
\textbf{Instructions.} Closed book. Write your answers on this paper, in the space provided. A
calculator is not needed --- the arithmetic is deliberately short. Answer every question --- there
is no negative marking. In Part C a reason is \textbf{optional}, but a correct reason can rescue a
wrong answer. In Parts D and E the marks are for the reasoning, not the vocabulary; an answer that
states what it cannot determine earns credit, one that asserts without a reason does not.\\[3pt]
\textbf{About the set letter.} Sets are handed out \textbf{at random}. Every set covers the same
topics at the same level, for the same marks, against the same marking scheme.}}
\vspace{3pt}


\parthead{A --- Multiple choice}{5 questions $\times$ 1 = 5 marks. Circle one letter.}

\begin{qA}
  \item \texttt{sum(1 for d in range(1, 7) if d \% 2 == 0) / 6} evaluates to\\
    (a) 0.5 \quad (b) 0.3333\\
    (c) 0.1667 \quad (d) 0.6667
  \item After \texttt{ws = [w for w in range(10) if w \% 3 == 0]}, \texttt{len(ws)} is\\
    (a) 3 \quad (b) 5\\
    (c) 4 \quad (d) an error
  \item \texttt{f = lambda w: (w - 3)**2}. The function returning $f'(w)$ is\\
    (a) \texttt{lambda w: (w - 3)**2} \quad (b) \texttt{lambda w: 2*(w - 3)}\\
    (c) \texttt{lambda w: 2*w - 3} \quad (d) \texttt{lambda w: w - 3}
  \item \texttt{np.mean(y)} rather than \texttt{np.median(y)} is the constant that minimises\\
    (a) \texttt{np.mean(np.abs(y - c))} \quad (b) \texttt{np.max(np.abs(y - c))}\\
    (c) neither of them \quad (d) \texttt{np.mean((y - c)**2)}
  \item In \texttt{w = w - eta * 2.0 * (w - 3.0)} run in a loop, the curvature is $2.0$, so the loop diverges as soon as \texttt{eta} reaches\\
    (a) \texttt{2.0} \quad (b) \texttt{0.5}\\
    (c) \texttt{1.0} \quad (d) no value --- it always converges
\end{qA}

\vspace{4pt}

\parthead{B --- Fill in the blanks}{4 $\times$ 1 = 4 marks. Write the missing word on the rule.}

\begin{qB}

  \item \texttt{grad(w)} returns \texttt{0.0} when \texttt{w} is at the minimum, because the derivative there equals \BL.

  \item \texttt{np.mean(y)} minimises \texttt{np.mean((y - c)**2)}; the constant minimising \texttt{np.mean(np.abs(y - c))} is \texttt{np.\BL(y)}.

  \item For a loop \texttt{w = w - eta * grad(w)} with largest curvature \texttt{L}, the loop converges only while \texttt{eta} is smaller than \BL.

  \item In \texttt{v = beta * v + g} followed by \texttt{w = w - eta * v}, the effective step size is \texttt{eta} divided by \BL.

\end{qB}

\vfill
\begin{center}\small
\begin{tabular}{|>{\centering\arraybackslash}m{1.9cm}|*{5}{>{\centering\arraybackslash}m{1.6cm}|}}
\hline
\rule{0pt}{12pt}\textbf{Question} & A1 & A2 & A3 & A4 & A5\\[2pt]\hline
\rule{0pt}{20pt}\textbf{Answer} & & & & & \\[7pt]\hline
\end{tabular}\\[2pt]
{\footnotesize Copy your circled Part A letters into this grid. The grid is what is marked.}
\end{center}

\newpage

\parthead{C --- True or false}{3 $\times$ 1 = 3 marks. Circle one. A reason is optional, but write one --- it can rescue a wrong answer.}

\begin{qC}

  \item A single call to \texttt{train\_test\_split} on a small dataset gives a reliable estimate of performance on new data. \hfill \textbf{TRUE} \; / \; \textbf{FALSE}\\[1pt]
    {\small Reason:}\,\blk{140mm}\rules{1}{6.4mm}

  \item Replacing \texttt{eta = 0.1} with \texttt{eta = 0.2} must always halve the number of loop iterations needed. \hfill \textbf{TRUE} \; / \; \textbf{FALSE}\\[1pt]
    {\small Reason:}\,\blk{140mm}\rules{1}{6.4mm}

  \item Setting \texttt{beta = 0.9} instead of \texttt{beta = 0.0} can make a well-conditioned problem take \emph{more} iterations, not fewer. \hfill \textbf{TRUE} \; / \; \textbf{FALSE}\\[1pt]
    {\small Reason:}\,\blk{140mm}\rules{1}{6.4mm}

\end{qC}

\vspace{4pt}

\parthead{D --- Short answer}{2 questions $\times$ 2 = 4 marks. Show how you decide, not just what you decide.}

\begin{qD}

  \item The loop below is run with \texttt{m = 5}, \texttt{eta = 0.3}.
\begin{code}
w = 0.0
for k in range(2):
    g = 2.0 * (w - m)
    w = w - eta * g
    print(k + 1, w)
\end{code}
    \begin{enumerate}[label=(\alph*),leftmargin=7mm,itemsep=1pt,topsep=2pt]
      \item Compute $w_1$ and $w_2$, showing your working.
      \item State the error multiplier $\lvert 1 - \eta f''(w) \rvert$, and say whether this run converges monotonically, converges while oscillating, or diverges.
    \end{enumerate}\par\vspace{68mm}

\end{qD}

\newpage

\parthead{D --- Short answer, continued}{}

\begin{qD}[start=2]

  \item \texttt{model.score} on the training rows improves at every epoch while \texttt{model.score} on the held-out rows falls at every epoch after the third. \textbf{(a)} Name what is happening. \textbf{(b)} State what you would change first, and why.\par\vspace{80mm}

\end{qD}

\vspace{4pt}

\parthead{E --- Reasoning}{1 question $\times$ 4 = 4 marks. You are marked on the reasoning you make visible.}

\begin{qE}

  \item The function below returns a fitted weight vector. Called with \texttt{StandardScaler().fit\_transform(X)}
and \texttt{eta=0.05} it converges in a few hundred iterations. Called with the \emph{raw} \texttt{X} --- one
column of which is an amount in rupees running into the millions --- and the same \texttt{eta=0.05}, it
returns \texttt{nan} within three iterations.

\begin{code}
def fit(X, y, eta=0.05, iters=500):
    w = np.zeros(X.shape[1])
    for _ in range(iters):
        g = 2.0 / len(y) * X.T @ (X @ w - y)
        w = w - eta * g
    return w
\end{code}


    \vspace{2mm}
    Answer all three parts overleaf. Any diagnosis that fits the evidence and is argued for earns
    full credit --- naming the ``expected'' fault is not what is marked.
    \begin{enumerate}[label=(\roman*),leftmargin=7mm,itemsep=2pt,topsep=2pt]
      \item Name the most likely cause, and explain \emph{why} it makes the run diverge when the same
            $\eta$ was stable before. \hfill \textbf{[2]}
      \item State one concrete check that would confirm or rule out your diagnosis. \hfill \textbf{[1]}
      \item Say what you would expect to see if your diagnosis turned out to be wrong. \hfill \textbf{[1]}
    \end{enumerate}

\end{qE}

\newpage

\parthead{E --- Reasoning, answer space}{Label each part (i), (ii), (iii). A sketch, a table or a list is as acceptable as prose.}

\vfill

{\footnotesize\itshape Continue on a continuation sheet if you need more room; label anything you write there with its question number.}

\newpage

\renewcommand{\SetLetter}{F}\setcounter{page}{1}


\begin{center}
{\large\bfseries CSAI2017P --- Applied Machine Learning (Theory)}\\[2pt]
{\large\bfseries Theory Quiz 2 (T02) \quad\textbullet\quad Set F}\\[3pt]
{\small Duration \textbf{40 minutes} \quad\textbullet\quad Maximum marks \textbf{20}
 \quad\textbullet\quad 15 questions on four pages}
\end{center}
\vspace{2pt}

% Column widths differ slightly from LQ01's: at LQ01's widths the row runs
% ~11.6pt past the margin and "Signature" does not fit its own cell.
\noindent\begin{tabular}{@{}p{1.4cm}p{5.0cm}p{1.9cm}p{3.5cm}p{1.3cm}p{2.4cm}@{}}
\textbf{Name} & \hrulefill & \textbf{SAP ID} & \hrulefill & \textbf{Batch} & \hrulefill\\[9pt]
\textbf{Date} & \hrulefill & \textbf{Signature} & \hrulefill & \textbf{Set} & \fbox{\rule{0pt}{11pt}\ F\ }\\
\end{tabular}

\vspace{4pt}
\noindent\fbox{\parbox{\dimexpr\linewidth-2\fboxsep-2\fboxrule}{\small
\textbf{Syllabus.} Prerequisite material --- probability, Python and elementary calculus --- together with
Unit I (Lecture 1) and Unit II (Lecture 2) as retained material: generalisation, cross-validation, and which
constant each loss fits; and, principally, \textbf{Unit III (Lecture 3)}: gradient descent and its update
rule, the learning rate and the stability threshold $\eta < 2/L$, batch, stochastic and mini-batch gradient
descent, and momentum.\\[3pt]
\textbf{Instructions.} Closed book. Write your answers on this paper, in the space provided. A
calculator is not needed --- the arithmetic is deliberately short. Answer every question --- there
is no negative marking. In Part C a reason is \textbf{optional}, but a correct reason can rescue a
wrong answer. In Parts D and E the marks are for the reasoning, not the vocabulary; an answer that
states what it cannot determine earns credit, one that asserts without a reason does not.\\[3pt]
\textbf{About the set letter.} Sets are handed out \textbf{at random}. Every set covers the same
topics at the same level, for the same marks, against the same marking scheme.}}
\vspace{3pt}


\parthead{A --- Multiple choice}{5 questions $\times$ 1 = 5 marks. Circle one letter.}

\begin{qA}
  \item An ordinary die is rolled once. The chance that the number showing is even is\\
    (a) one in three \quad (b) one in two\\
    (c) one in six \quad (d) two in three
  \item Counting the whole numbers from zero to nine that divide exactly by three, the count is\\
    (a) 3 \quad (b) 5\\
    (c) none of them do \quad (d) 4
  \item A quantity grows as the square of how far a setting lies from three. The rate at which it changes as the setting moves is\\
    (a) the square of the distance from three \quad (b) twice the setting, less three\\
    (c) twice the distance from three \quad (d) the distance from three
  \item One number is to stand in for a set of readings, chosen so that the total of the \emph{squared} differences is as small as possible. That number is their\\
    (a) mean \quad (b) median\\
    (c) mode \quad (d) largest value
  \item A step size that is small enough works, and one that is a little too large fails outright rather than merely working less well. The boundary between the two is\\
    (a) always one \quad (b) half the largest curvature\\
    (c) there is no boundary \quad (d) twice the reciprocal of the largest curvature
\end{qA}

\vspace{4pt}

\parthead{B --- Fill in the blanks}{4 $\times$ 1 = 4 marks. Write the missing word on the rule.}

\begin{qB}

  \item Where a smooth curve is at its lowest, its slope is neither rising nor falling: the derivative there equals \BL.

  \item Squaring the errors before averaging makes the best single value the mean; taking their size instead makes it the \BL.

  \item Writing $L$ for the largest curvature of the surface, the method converges only while the step size is smaller than \BL.

  \item Keeping nine tenths of the previous velocity makes the steps about ten times as long, because the effective step is the step size divided by \BL.

\end{qB}

\vfill
\begin{center}\small
\begin{tabular}{|>{\centering\arraybackslash}m{1.9cm}|*{5}{>{\centering\arraybackslash}m{1.6cm}|}}
\hline
\rule{0pt}{12pt}\textbf{Question} & A1 & A2 & A3 & A4 & A5\\[2pt]\hline
\rule{0pt}{20pt}\textbf{Answer} & & & & & \\[7pt]\hline
\end{tabular}\\[2pt]
{\footnotesize Copy your circled Part A letters into this grid. The grid is what is marked.}
\end{center}

\newpage

\parthead{C --- True or false}{3 $\times$ 1 = 3 marks. Circle one. A reason is optional, but write one --- it can rescue a wrong answer.}

\begin{qC}

  \item Splitting a small dataset in two once, and quoting the score on the second half, gives a reliable estimate of performance on new data. \hfill \textbf{TRUE} \; / \; \textbf{FALSE}\\[1pt]
    {\small Reason:}\,\blk{140mm}\rules{1}{6.4mm}

  \item Taking steps twice as large always halves the number of steps needed to arrive. \hfill \textbf{TRUE} \; / \; \textbf{FALSE}\\[1pt]
    {\small Reason:}\,\blk{140mm}\rules{1}{6.4mm}

  \item Carrying momentum can make a run take longer than it would without it, when the surface is an even bowl rather than a narrow valley. \hfill \textbf{TRUE} \; / \; \textbf{FALSE}\\[1pt]
    {\small Reason:}\,\blk{140mm}\rules{1}{6.4mm}

\end{qC}

\vspace{4pt}

\parthead{D --- Short answer}{2 questions $\times$ 2 = 4 marks. Show how you decide, not just what you decide.}

\begin{qD}

  \item A quantity is smallest when a setting equals ten, and grows as the square of the distance from ten. Starting the setting at zero and repeatedly subtracting the step size times the slope, with a step size of one tenth,
    \begin{enumerate}[label=(\alph*),leftmargin=7mm,itemsep=1pt,topsep=2pt]
      \item Compute $w_1$ and $w_2$, showing your working.
      \item State the error multiplier $\lvert 1 - \eta f''(w) \rvert$, and say whether this run converges monotonically, converges while oscillating, or diverges.
    \end{enumerate}\par\vspace{68mm}

\end{qD}

\newpage

\parthead{D --- Short answer, continued}{}

\begin{qD}[start=2]

  \item As training goes on, the model does steadily better on the examples it is being fitted to, and steadily worse on the examples held back from it. \textbf{(a)} Name what is happening. \textbf{(b)} State what you would change first, and why.\par\vspace{80mm}

\end{qD}

\vspace{4pt}

\parthead{E --- Reasoning}{1 question $\times$ 4 = 4 marks. You are marked on the reasoning you make visible.}

\begin{qE}

  \item A team fits a model by repeatedly stepping downhill. On features that have each been centred and
divided by their spread, with a step size of five hundredths, the fit settles after a few hundred steps. The
team then rebuilds the same thing on the original untouched columns --- one of which is an amount of money
running into the millions --- keeps the step size exactly as it was, and within three steps the loss is no
longer a number at all.


    \vspace{2mm}
    Answer all three parts overleaf. Any diagnosis that fits the evidence and is argued for earns
    full credit --- naming the ``expected'' fault is not what is marked.
    \begin{enumerate}[label=(\roman*),leftmargin=7mm,itemsep=2pt,topsep=2pt]
      \item Name the most likely cause, and explain \emph{why} it makes the run diverge when the same
            $\eta$ was stable before. \hfill \textbf{[2]}
      \item State one concrete check that would confirm or rule out your diagnosis. \hfill \textbf{[1]}
      \item Say what you would expect to see if your diagnosis turned out to be wrong. \hfill \textbf{[1]}
    \end{enumerate}

\end{qE}

\newpage

\parthead{E --- Reasoning, answer space}{Label each part (i), (ii), (iii). A sketch, a table or a list is as acceptable as prose.}

\vfill

{\footnotesize\itshape Continue on a continuation sheet if you need more room; label anything you write there with its question number.}

\newpage

\renewcommand{\SetLetter}{G}\setcounter{page}{1}


\begin{center}
{\large\bfseries CSAI2017P --- Applied Machine Learning (Theory)}\\[2pt]
{\large\bfseries Theory Quiz 2 (T02) \quad\textbullet\quad Set G}\\[3pt]
{\small Duration \textbf{40 minutes} \quad\textbullet\quad Maximum marks \textbf{20}
 \quad\textbullet\quad 15 questions on four pages}
\end{center}
\vspace{2pt}

% Column widths differ slightly from LQ01's: at LQ01's widths the row runs
% ~11.6pt past the margin and "Signature" does not fit its own cell.
\noindent\begin{tabular}{@{}p{1.4cm}p{5.0cm}p{1.9cm}p{3.5cm}p{1.3cm}p{2.4cm}@{}}
\textbf{Name} & \hrulefill & \textbf{SAP ID} & \hrulefill & \textbf{Batch} & \hrulefill\\[9pt]
\textbf{Date} & \hrulefill & \textbf{Signature} & \hrulefill & \textbf{Set} & \fbox{\rule{0pt}{11pt}\ G\ }\\
\end{tabular}

\vspace{4pt}
\noindent\fbox{\parbox{\dimexpr\linewidth-2\fboxsep-2\fboxrule}{\small
\textbf{Syllabus.} Prerequisite material --- probability, Python and elementary calculus --- together with
Unit I (Lecture 1) and Unit II (Lecture 2) as retained material: generalisation, cross-validation, and which
constant each loss fits; and, principally, \textbf{Unit III (Lecture 3)}: gradient descent and its update
rule, the learning rate and the stability threshold $\eta < 2/L$, batch, stochastic and mini-batch gradient
descent, and momentum.\\[3pt]
\textbf{Instructions.} Closed book. Write your answers on this paper, in the space provided. A
calculator is not needed --- the arithmetic is deliberately short. Answer every question --- there
is no negative marking. In Part C a reason is \textbf{optional}, but a correct reason can rescue a
wrong answer. In Parts D and E the marks are for the reasoning, not the vocabulary; an answer that
states what it cannot determine earns credit, one that asserts without a reason does not.\\[3pt]
\textbf{About the set letter.} Sets are handed out \textbf{at random}. Every set covers the same
topics at the same level, for the same marks, against the same marking scheme.}}
\vspace{3pt}


\parthead{A --- Multiple choice}{5 questions $\times$ 1 = 5 marks. Circle one letter.}

\begin{qA}
  \item Compare, for one roll of a fair die, $P(\text{even})$ with $P(\text{a six})$\\
    (a) the second is larger \quad (b) they are equal\\
    (c) the first is three times the second \quad (d) cannot be decided
  \item Compare \texttt{len([w for w in range(10) if w \% 3 == 0])} with \texttt{len(range(3))}\\
    (a) the first is larger \quad (b) the second is larger\\
    (c) they are equal \quad (d) one raises an error
  \item Compare $f(w) = (w-3)^2$ with its derivative $f'(w)$ at $w = 5$\\
    (a) $f = 4$ and $f' = 2$ \quad (b) $f = 2$ and $f' = 4$\\
    (c) $f = 4$ and $f' = 7$ \quad (d) $f = 4$ and $f' = 4$
  \item Compare the constant minimising the squared error with the constant minimising the absolute error\\
    (a) the first is the median, the second the mean \quad (b) the first is the mean, the second the median\\
    (c) both are the mean \quad (d) both are the median
  \item Compare $\eta = 0.10$ with $\eta = 0.90$ on $f(w) = (w-3)^2$, where both give $\lvert 1 - 2\eta \rvert = 0.80$\\
    (a) they converge at the same rate, one oscillating \quad (b) the second is nine times faster\\
    (c) the second diverges \quad (d) the first diverges
\end{qA}

\vspace{4pt}

\parthead{B --- Fill in the blanks}{4 $\times$ 1 = 4 marks. Write the missing word on the rule.}

\begin{qB}

  \item Compare the derivative away from a minimum with the derivative at one: the second equals \BL.

  \item Compare the two: squared error is fitted by the mean, absolute error by the \BL.

  \item Compare a converging run with a diverging one: the boundary between them is $\eta$ equal to \BL.

  \item Compare plain gradient descent with momentum: the second takes an effective step of $\eta$ divided by \BL.

\end{qB}

\vfill
\begin{center}\small
\begin{tabular}{|>{\centering\arraybackslash}m{1.9cm}|*{5}{>{\centering\arraybackslash}m{1.6cm}|}}
\hline
\rule{0pt}{12pt}\textbf{Question} & A1 & A2 & A3 & A4 & A5\\[2pt]\hline
\rule{0pt}{20pt}\textbf{Answer} & & & & & \\[7pt]\hline
\end{tabular}\\[2pt]
{\footnotesize Copy your circled Part A letters into this grid. The grid is what is marked.}
\end{center}

\newpage

\parthead{C --- True or false}{3 $\times$ 1 = 3 marks. Circle one. A reason is optional, but write one --- it can rescue a wrong answer.}

\begin{qC}

  \item Between one train/test split and five-fold cross-validation on a small dataset, the first is just as reliable an estimate of performance on new data. \hfill \textbf{TRUE} \; / \; \textbf{FALSE}\\[1pt]
    {\small Reason:}\,\blk{140mm}\rules{1}{6.4mm}

  \item Between a run at $\eta$ and the same run at $2\eta$, the second must always take half as many iterations. \hfill \textbf{TRUE} \; / \; \textbf{FALSE}\\[1pt]
    {\small Reason:}\,\blk{140mm}\rules{1}{6.4mm}

  \item Between plain gradient descent and momentum on a well-conditioned problem, the second can be the one that needs more iterations. \hfill \textbf{TRUE} \; / \; \textbf{FALSE}\\[1pt]
    {\small Reason:}\,\blk{140mm}\rules{1}{6.4mm}

\end{qC}

\vspace{4pt}

\parthead{D --- Short answer}{2 questions $\times$ 2 = 4 marks. Show how you decide, not just what you decide.}

\begin{qD}

  \item Compare two runs on $f(w) = (w-10)^2$ from $w_0 = 0$, one with $\eta = 0.7$ and one with $\eta = 0.05$. Take the \textbf{first} of them, $\eta = 0.7$, with $w \leftarrow w - \eta f'(w)$.
    \begin{enumerate}[label=(\alph*),leftmargin=7mm,itemsep=1pt,topsep=2pt]
      \item Compute $w_1$ and $w_2$, showing your working.
      \item State the error multiplier $\lvert 1 - \eta f''(w) \rvert$, and say whether this run converges monotonically, converges while oscillating, or diverges.
    \end{enumerate}\par\vspace{68mm}

\end{qD}

\newpage

\parthead{D --- Short answer, continued}{}

\begin{qD}[start=2]

  \item Compare the two curves recorded for one training run: the error on the rows the model was fitted to falls steadily, while the error on the rows held back rises steadily. \textbf{(a)} Name what is happening. \textbf{(b)} State what you would change first, and why.\par\vspace{80mm}

\end{qD}

\vspace{4pt}

\parthead{E --- Reasoning}{1 question $\times$ 4 = 4 marks. You are marked on the reasoning you make visible.}

\begin{qE}

  \item Compare two runs of the same linear model, fitted by gradient descent with $\eta = 0.05$ in both. The
first uses \textbf{standardised} features and converges in a few hundred iterations. The second uses the
\textbf{raw} features --- one of which is an amount in rupees running into the millions --- changes nothing
else, and reaches a loss of \texttt{nan} within three iterations.


    \vspace{2mm}
    Answer all three parts overleaf. Any diagnosis that fits the evidence and is argued for earns
    full credit --- naming the ``expected'' fault is not what is marked.
    \begin{enumerate}[label=(\roman*),leftmargin=7mm,itemsep=2pt,topsep=2pt]
      \item Name the most likely cause, and explain \emph{why} it makes the run diverge when the same
            $\eta$ was stable before. \hfill \textbf{[2]}
      \item State one concrete check that would confirm or rule out your diagnosis. \hfill \textbf{[1]}
      \item Say what you would expect to see if your diagnosis turned out to be wrong. \hfill \textbf{[1]}
    \end{enumerate}

\end{qE}

\newpage

\parthead{E --- Reasoning, answer space}{Label each part (i), (ii), (iii). A sketch, a table or a list is as acceptable as prose.}

\vfill

{\footnotesize\itshape Continue on a continuation sheet if you need more room; label anything you write there with its question number.}

\newpage

\renewcommand{\SetLetter}{H}\setcounter{page}{1}


\begin{center}
{\large\bfseries CSAI2017P --- Applied Machine Learning (Theory)}\\[2pt]
{\large\bfseries Theory Quiz 2 (T02) \quad\textbullet\quad Set H}\\[3pt]
{\small Duration \textbf{40 minutes} \quad\textbullet\quad Maximum marks \textbf{20}
 \quad\textbullet\quad 15 questions on four pages}
\end{center}
\vspace{2pt}

% Column widths differ slightly from LQ01's: at LQ01's widths the row runs
% ~11.6pt past the margin and "Signature" does not fit its own cell.
\noindent\begin{tabular}{@{}p{1.4cm}p{5.0cm}p{1.9cm}p{3.5cm}p{1.3cm}p{2.4cm}@{}}
\textbf{Name} & \hrulefill & \textbf{SAP ID} & \hrulefill & \textbf{Batch} & \hrulefill\\[9pt]
\textbf{Date} & \hrulefill & \textbf{Signature} & \hrulefill & \textbf{Set} & \fbox{\rule{0pt}{11pt}\ H\ }\\
\end{tabular}

\vspace{4pt}
\noindent\fbox{\parbox{\dimexpr\linewidth-2\fboxsep-2\fboxrule}{\small
\textbf{Syllabus.} Prerequisite material --- probability, Python and elementary calculus --- together with
Unit I (Lecture 1) and Unit II (Lecture 2) as retained material: generalisation, cross-validation, and which
constant each loss fits; and, principally, \textbf{Unit III (Lecture 3)}: gradient descent and its update
rule, the learning rate and the stability threshold $\eta < 2/L$, batch, stochastic and mini-batch gradient
descent, and momentum.\\[3pt]
\textbf{Instructions.} Closed book. Write your answers on this paper, in the space provided. A
calculator is not needed --- the arithmetic is deliberately short. Answer every question --- there
is no negative marking. In Part C a reason is \textbf{optional}, but a correct reason can rescue a
wrong answer. In Parts D and E the marks are for the reasoning, not the vocabulary; an answer that
states what it cannot determine earns credit, one that asserts without a reason does not.\\[3pt]
\textbf{About the set letter.} Sets are handed out \textbf{at random}. Every set covers the same
topics at the same level, for the same marks, against the same marking scheme.}}
\vspace{3pt}


\parthead{A --- Multiple choice}{5 questions $\times$ 1 = 5 marks. Circle one letter.}

\begin{qA}
  \item \emph{For one toss of a fair coin, $P(\text{head}) = 1/2$.} For one roll of a fair die, $P(\text{even})$ is\\
    (a) $1/3$ \quad (b) $1/6$\\
    (c) $2/3$ \quad (d) $1/2$
  \item \emph{\texttt{len([w for w in range(6) if w \% 2 == 0])} is $3$.} Then \texttt{len([w for w in range(10) if w \% 3 == 0])} is\\
    (a) 3 \quad (b) 4\\
    (c) 5 \quad (d) an error
  \item \emph{The derivative of $(w-1)^2$ is $2(w-1)$.} The derivative of $(w-3)^2$ is\\
    (a) $2(w-3)$ \quad (b) $(w-3)^2$\\
    (c) $2w-3$ \quad (d) $w-3$
  \item \emph{The constant minimising $\sum_i \lvert y_i - c\rvert$ is the median.} The constant minimising $\sum_i (y_i - c)^2$ is the\\
    (a) median \quad (b) mode\\
    (c) mean \quad (d) largest value
  \item \emph{For $f(w) = 3w^2$ the curvature is $6$, so the run converges only while $\eta < 1/3$.} For $f(w) = (w-3)^2$ the curvature is $2$, so it converges only while\\
    (a) $\eta < 2$ \quad (b) $\eta < 1$\\
    (c) $\eta < 0.5$ \quad (d) any $\eta > 0$
\end{qA}

\vspace{4pt}

\parthead{B --- Fill in the blanks}{4 $\times$ 1 = 4 marks. Write the missing word on the rule.}

\begin{qB}

  \item \emph{Away from the minimum of $(w-3)^2$ the derivative is non-zero.} At the minimum itself it equals \BL.

  \item \emph{Squared error is minimised by the mean.} Absolute error is minimised by the \BL.

  \item \emph{For curvature $L = 2$ the threshold is $\eta < 1$.} In general the run converges only while $\eta$ is smaller than \BL.

  \item \emph{With $\beta = 0.99$ the effective step is $100\eta$.} In general it is $\eta$ divided by \BL.

\end{qB}

\vfill
\begin{center}\small
\begin{tabular}{|>{\centering\arraybackslash}m{1.9cm}|*{5}{>{\centering\arraybackslash}m{1.6cm}|}}
\hline
\rule{0pt}{12pt}\textbf{Question} & A1 & A2 & A3 & A4 & A5\\[2pt]\hline
\rule{0pt}{20pt}\textbf{Answer} & & & & & \\[7pt]\hline
\end{tabular}\\[2pt]
{\footnotesize Copy your circled Part A letters into this grid. The grid is what is marked.}
\end{center}

\newpage

\parthead{C --- True or false}{3 $\times$ 1 = 3 marks. Circle one. A reason is optional, but write one --- it can rescue a wrong answer.}

\begin{qC}

  \item \emph{Five-fold cross-validation quotes a mean and a spread.} A single train/test split on a small dataset is therefore just as reliable an estimate. \hfill \textbf{TRUE} \; / \; \textbf{FALSE}\\[1pt]
    {\small Reason:}\,\blk{140mm}\rules{1}{6.4mm}

  \item \emph{Going from $\eta = 0.1$ to $\eta = 0.5$ on $f(w) = (w-3)^2$ takes the run from many steps to exactly one.} Therefore doubling any learning rate always halves the iterations needed. \hfill \textbf{TRUE} \; / \; \textbf{FALSE}\\[1pt]
    {\small Reason:}\,\blk{140mm}\rules{1}{6.4mm}

  \item \emph{In a ravine, momentum cut $954$ iterations to $135$.} On a well-conditioned problem it can instead need more iterations than plain gradient descent. \hfill \textbf{TRUE} \; / \; \textbf{FALSE}\\[1pt]
    {\small Reason:}\,\blk{140mm}\rules{1}{6.4mm}

\end{qC}

\vspace{4pt}

\parthead{D --- Short answer}{2 questions $\times$ 2 = 4 marks. Show how you decide, not just what you decide.}

\begin{qD}

  \item \emph{For $f(w) = (w-2)^2$ from $w_0 = 0$ with $\eta = 0.5$, the multiplier $\lvert 1 - 2\eta \rvert$ is $0$ and one step lands exactly on $w = 2$.} Now take $f(w) = (w-6)^2$ from $w_0 = 0$ with $\eta = 0.9$ and $w \leftarrow w - \eta f'(w)$.
    \begin{enumerate}[label=(\alph*),leftmargin=7mm,itemsep=1pt,topsep=2pt]
      \item Compute $w_1$ and $w_2$, showing your working.
      \item State the error multiplier $\lvert 1 - \eta f''(w) \rvert$, and say whether this run converges monotonically, converges while oscillating, or diverges.
    \end{enumerate}\par\vspace{68mm}

\end{qD}

\newpage

\parthead{D --- Short answer, continued}{}

\begin{qD}[start=2]

  \item \emph{A model whose training error and held-out error fall together is generalising.} A model whose training loss falls smoothly towards zero while its held-out error rises steadily is doing something else. \textbf{(a)} Name it. \textbf{(b)} State what you would change first, and why.\par\vspace{80mm}

\end{qD}

\vspace{4pt}

\parthead{E --- Reasoning}{1 question $\times$ 4 = 4 marks. You are marked on the reasoning you make visible.}

\begin{qE}

  \item \emph{Worked analogue. A run that was stable at $\eta = 0.01$ is switched to $\beta = 0.99$ with $\eta$
unchanged, and diverges: the effective step became $100\eta$ rather than $10\eta$. The check is to divide
$\eta$ by ten and re-run; if it is now stable, the diagnosis holds.}

\vspace{2mm}
Now consider this case. A team fits a linear model by gradient descent on \textbf{standardised} features with
$\eta = 0.05$, and it converges in a few hundred iterations. They rebuild the same pipeline on the
\textbf{raw} features --- one of which is an amount in rupees running into the millions --- keep
$\eta = 0.05$, and the loss becomes \texttt{nan} within three iterations.


    \vspace{2mm}
    Answer all three parts overleaf. Any diagnosis that fits the evidence and is argued for earns
    full credit --- naming the ``expected'' fault is not what is marked.
    \begin{enumerate}[label=(\roman*),leftmargin=7mm,itemsep=2pt,topsep=2pt]
      \item Name the most likely cause, and explain \emph{why} it makes the run diverge when the same
            $\eta$ was stable before. \hfill \textbf{[2]}
      \item State one concrete check that would confirm or rule out your diagnosis. \hfill \textbf{[1]}
      \item Say what you would expect to see if your diagnosis turned out to be wrong. \hfill \textbf{[1]}
    \end{enumerate}

\end{qE}

\newpage

\parthead{E --- Reasoning, answer space}{Label each part (i), (ii), (iii). A sketch, a table or a list is as acceptable as prose.}

\vfill

{\footnotesize\itshape Continue on a continuation sheet if you need more room; label anything you write there with its question number.}

\newpage

\end{document}